\documentclass{article}
\usepackage{amsmath,amssymb,amsthm}
\usepackage{algorithm,algorithmic}
\usepackage{bm}
\usepackage{hyperref}
\newtheorem{theorem}{Theorem}
\newtheorem{lemma}{Lemma}
\newtheorem{assumption}{Assumption}
\newcommand{\gap}{\operatorname{Gap}}
\newcommand{\E}{\mathbb{E}}
\newcommand{\R}{\mathbb{R}}

\begin{document}

\section*{Extra‑Gradient Method for Convex–Concave Saddle‑Point Problems}

\section*{Mathematical Formulation}
Consider a continuously differentiable function 
\[
L:\mathcal{X}\times\mathcal{Y}\to\R,
\qquad \mathcal{X}\subset\R^{n},\;\mathcal{Y}\subset\R^{m}
\]
which is convex in the first argument and concave in the second argument.
Define the monotone operator
\[
F(z)\;:=\;
\begin{bmatrix}
\nabla_{x}L(x,y)\\[2mm]
-\nabla_{y}L(x,y)
\end{bmatrix},
\qquad 
z:=\begin{bmatrix}x\\y\end{bmatrix}\in\mathcal{Z}:=\mathcal{X}\times\mathcal{Y}.
\]

Given a stepsize $\eta>0$, the (deterministic) Extra‑Gradient (EG) iteration is
\[
\boxed{
\begin{aligned}
z_{k+\frac12} &:= \Pi_{\mathcal{Z}}\!\bigl(z_{k}-\eta\,F(z_{k})\bigr),\\[1mm]
z_{k+1}       &:= \Pi_{\mathcal{Z}}\!\bigl(z_{k}-\eta\,F(z_{k+\frac12})\bigr),
\end{aligned}}
\tag{EG}
\]
where $\Pi_{\mathcal{Z}}$ denotes Euclidean projection onto $\mathcal{Z}$.
Writing the components explicitly,
\[
\begin{aligned}
x_{k+\frac12}&= \Pi_{\mathcal{X}}\!\bigl(x_{k}-\eta\,\nabla_{x}L(x_{k},y_{k})\bigr),\\
y_{k+\frac12}&= \Pi_{\mathcal{Y}}\!\bigl(y_{k}+\eta\,\nabla_{y}L(x_{k},y_{k})\bigr),\\[1mm]
x_{k+1}&= \Pi_{\mathcal{X}}\!\bigl(x_{k}-\eta\,\nabla_{x}L(x_{k+\frac12},y_{k+\frac12})\bigr),\\
y_{k+1}&= \Pi_{\mathcal{Y}}\!\bigl(y_{k}+\eta\,\nabla_{y}L(x_{k+\frac12},y_{k+\frac12})\bigr).
\end{aligned}
\]

\section*{Pseudocode}
\begin{algorithm}[H]
\caption{Extra‑Gradient for $\displaystyle\min_{x\in\mathcal{X}}\max_{y\in\mathcal{Y}}L(x,y)$}
\begin{algorithmic}[1]
\REQUIRE Initial point $z_{0}=(x_{0},y_{0})\in\mathcal{Z}$, stepsize $\eta>0$, 
maximum iterations $K$.
\FOR{$k=0,1,\dots,K-1$}
    \STATE $g_{k}=F(z_{k})$                                 \COMMENT{gradient at current iterate}
    \STATE $z_{k+\frac12}= \Pi_{\mathcal{Z}}\!\bigl(z_{k}-\eta g_{k}\bigr)$
    \STATE $\tilde g_{k}=F(z_{k+\frac12})$                    \COMMENT{gradient at the extrapolated point}
    \STATE $z_{k+1}= \Pi_{\mathcal{Z}}\!\bigl(z_{k}-\eta \tilde g_{k}\bigr)$
\ENDFOR
\RETURN $z_{K}$
\end{algorithmic}
\end{algorithm}

\section*{Dry Run Example}
We illustrate the EG update on a simple **pure minimisation** problem
\[
f(w)=(w-3)^{2},\qquad w\in\R .
\]
Here $L(x,y)=f(x)$, $\mathcal{Y}$ is empty, and $F(w)=\nabla f(w)=2(w-3)$.  
Take $w_{0}=0$, stepsize $\eta=0.1$, and run three iterations.

\begin{center}
\begin{tabular}{c|c|c|c}
\textbf{Iter.} & $w_{k}$ & $w_{k+\frac12}=w_{k}-\eta\nabla f(w_{k})$ & $w_{k+1}=w_{k}-\eta\nabla f(w_{k+\frac12})$\\\hline
0 & $0$ & $0-0.1\cdot2(0-3)=0+0.6=0.6$ & $0-0.1\cdot2(0.6-3)=0-0.1\cdot(-4.8)=0.48$\\[2mm]
1 & $0.48$ & $0.48-0.1\cdot2(0.48-3)=0.48+0.504=0.984$ & $0.48-0.1\cdot2(0.984-3)=0.48-0.1\cdot(-4.032)=0.8832$\\[2mm]
2 & $0.8832$ & $0.8832-0.1\cdot2(0.8832-3)=0.8832+0.42336=1.30656$ &
$0.8832-0.1\cdot2(1.30656-3)=0.8832-0.1\cdot(-3.38688)=1.221888$
\end{tabular}
\end{center}

The objective values are
\[
f(w_{0})=9,\; f(w_{1})\approx(0.48-3)^{2}=6.3504,\;
f(w_{2})\approx(0.8832-3)^{2}=4.4777,\;
f(w_{3})\approx(1.2219-3)^{2}=3.162.
\]
The sequence moves toward the minimiser $w^{\star}=3$.

\section*{Assumptions}
\begin{assumption}[Convex–concave structure]\label{as:convconc}
For all $x\in\mathcal{X}$ the map $y\mapsto L(x,y)$ is concave, and for all $y\in\mathcal{Y}$ the map $x\mapsto L(x,y)$ is convex.
\end{assumption}
\begin{assumption}[Lipschitz gradient]\label{as:smooth}
There exists $L>0$ such that for all $z,z'\in\mathcal{Z}$,
\[
\|F(z)-F(z')\|\le L\|z-z'\|.
\tag{1}
\]
\end{assumption}
\begin{assumption}[Bounded feasible set]\label{as:bounded}
The domain $\mathcal{Z}$ is closed, convex and has diameter $D$,
\[
\max_{z,z'\in\mathcal{Z}}\|z-z'\|\le D.
\tag{2}
\]
\end{assumption}
\begin{assumption}[Existence of a saddle point]\label{as:saddle}
There exists $z^{\star}=(x^{\star},y^{\star})\in\mathcal{Z}$ such that
\[
L(x^{\star},y)\le L(x^{\star},y^{\star})\le L(x,y^{\star}),\qquad\forall x\in\mathcal{X},y\in\mathcal{Y}.
\tag{3}
\]
\end{assumption}

\section*{Main Convergence Theorem}
\begin{theorem}[Ergodic $O(1/T)$ rate]\label{thm:rate}
Let $\{z_{k}\}_{k\ge0}$ be generated by \eqref{EG} with a constant stepsize 
\[
\eta\le \frac{1}{2L}.
\tag{4}
\]
Define the averaged iterate
\[
\bar z_{T}:=\frac{1}{T}\sum_{k=0}^{T-1}z_{k+\frac12}.
\]
Then for any $z\in\mathcal{Z}$
\[
\gap(\bar z_{T};z):=
\max_{y\in\mathcal{Y}}L(\bar x_{T},y)-\min_{x\in\mathcal{X}}L(x,\bar y_{T})
\le \frac{D^{2}}{2\eta T},
\tag{5}
\]
where $D$ is the diameter defined in \eqref{as:bounded}. In particular,
\[
\gap(\bar z_{T})\le \frac{2L D^{2}}{T}.
\tag{6}
\]
\end{theorem}

\section*{Theoretical Properties}
\begin{itemize}
\item \textbf{Stability.} The projection step guarantees that all iterates stay inside the compact set $\mathcal{Z}$, preventing divergence even when $F$ is only monotone.
\item \textbf{Stepsize sensitivity.} Condition \eqref{thm:rate}\,(4) is tight for the standard analysis; any $\eta>1/(2L)$ may break the contraction used in the proof.
\item \textbf{Comparison.} Classical Gradient Descent–Ascent (GDA) with the same stepsize yields $O(1/\sqrt{T})$ in the worst case for bilinear games, whereas EG attains the optimal $O(1/T)$ rate under the same smoothness assumptions.
\end{itemize}

\section*{Preliminaries (Definitions and Lemmas)}
\begin{definition}[Monotone operator]
$F$ is monotone on $\mathcal{Z}$ if
\[
\langle F(z)-F(z'),\,z-z'\rangle\ge0,\qquad\forall z,z'\in\mathcal{Z}.
\tag{7}
\]
\end{definition}
\begin{lemma}[Lipschitz $\Rightarrow$ monotone for convex–concave $L$]\label{lem:mono}
If $L$ satisfies Assumptions~\ref{as:convconc} and \ref{as:smooth}, then $F$ defined above is monotone.
\end{lemma}
\begin{proof}
For any $z=(x,y)$ and $z'=(x',y')$,
\[
\begin{aligned}
&\langle F(z)-F(z'),z-z'\rangle\\
&=\langle\nabla_{x}L(x,y)-\nabla_{x}L(x',y'),x-x'\rangle
   -\langle\nabla_{y}L(x,y)-\nabla_{y}L(x',y'),y-y'\rangle.
\end{aligned}
\]
Convexity in $x$ implies 
$\langle\nabla_{x}L(x,y)-\nabla_{x}L(x',y),x-x'\rangle\ge0$,
concavity in $y$ implies 
$-\langle\nabla_{y}L(x,y)-\nabla_{y}L(x,y'),y-y'\rangle\ge0$.
Summing the two non‑negative terms yields \eqref{7}. ∎
\end{proof}

\section*{Main Proof (Step‑by‑Step Derivation)}
We prove Theorem~\ref{thm:rate}. Let $z^{\star}$ be any saddle point (exists by Assumption~\ref{as:saddle}).

\noindent\textbf{[Step 1]} \emph{One‑step inequality for the extrapolation.}  
Using the definition of the projection,
\[
\|z_{k+\frac12}-z^{\star}\|^{2}
\le\|z_{k}-\eta F(z_{k})-z^{\star}\|^{2}.
\tag{8}
\]
Expanding the right‑hand side,
\[
\begin{aligned}
\|z_{k}-\eta F(z_{k})-z^{\star}\|^{2}
&=\|z_{k}-z^{\star}\|^{2}
   -2\eta\langle F(z_{k}),z_{k}-z^{\star}\rangle
   +\eta^{2}\|F(z_{k})\|^{2}. \tag{9}
\end{aligned}
\]

\noindent\textbf{[Step 2]} \emph{Monotonicity bound.}  
By Lemma~\ref{lem:mono},
\[
\langle F(z_{k}),z_{k}-z^{\star}\rangle\ge0.
\tag{10}
\]
Hence the middle term in \eqref{9} is non‑negative and can be dropped for an upper bound:
\[
\|z_{k+\frac12}-z^{\star}\|^{2}
\le\|z_{k}-z^{\star}\|^{2}+\eta^{2}\|F(z_{k})\|^{2}.
\tag{11}
\]

\noindent\textbf{[Step 3]} \emph{One‑step inequality for the final update.}  
Similarly,
\[
\|z_{k+1}-z^{\star}\|^{2}
\le\|z_{k}-\eta F(z_{k+\frac12})-z^{\star}\|^{2}.
\tag{12}
\]
Expanding,
\[
\|z_{k+1}-z^{\star}\|^{2}
\le\|z_{k}-z^{\star}\|^{2}
   -2\eta\langle F(z_{k+\frac12}),z_{k}-z^{\star}\rangle
   +\eta^{2}\|F(z_{k+\frac12})\|^{2}.
\tag{13}
\]

\noindent\textbf{[Step 4]} \emph{Replace $z_{k}$ by $z_{k+\frac12}$ in the inner product.}  
Add and subtract $\langle F(z_{k+\frac12}),z_{k+\frac12}-z^{\star}\rangle$:
\[
\begin{aligned}
-2\eta\langle F(z_{k+\frac12}),z_{k}-z^{\star}\rangle
&=-2\eta\langle F(z_{k+\frac12}),z_{k+\frac12}-z^{\star}\rangle \\
&\quad-2\eta\langle F(z_{k+\frac12}),z_{k}-z_{k+\frac12}\rangle .
\end{aligned}
\tag{14}
\]

\noindent\textbf{[Step 5]} \emph{Use monotonicity again.}  
Since $F$ is monotone,
\[
\langle F(z_{k+\frac12})-F(z_{k}),z_{k+\frac12}-z_{k}\rangle\ge0
\;\Longrightarrow\;
-\langle F(z_{k+\frac12}),z_{k}-z_{k+\frac12}\rangle
\le -\langle F(z_{k}),z_{k}-z_{k+\frac12}\rangle .
\tag{15}
\]

\noindent\textbf{[Step 6]} \emph{Bound the cross term with Lipschitzness.}  
From \eqref{as:smooth} and Cauchy–Schwarz,
\[
\begin{aligned}
\bigl|\langle F(z_{k+\frac12})-F(z_{k}),z_{k}-z_{k+\frac12}\rangle\bigr|
&\le\|F(z_{k+\frac12})-F(z_{k})\|\;\|z_{k}-z_{k+\frac12}\|\\
&\le L\|z_{k+\frac12}-z_{k}\|^{2}.
\end{aligned}
\tag{16}
\]

\noindent\textbf{[Step 7]} \emph{Express $z_{k+\frac12}-z_{k}$ via the first update.}  
From the definition of $z_{k+\frac12}$,
\[
z_{k+\frac12}= \Pi_{\mathcal{Z}}(z_{k}-\eta F(z_{k}))
\;\Longrightarrow\;
\|z_{k+\frac12}-z_{k}\|\le\eta\|F(z_{k})\|.
\tag{17}
\]

\noindent\textbf{[Step 8]} \emph{Combine \eqref{13}–\eqref{17}.}  
Using \eqref{15}–\eqref{17} we obtain
\[
\begin{aligned}
\|z_{k+1}-z^{\star}\|^{2}
&\le\|z_{k}-z^{\star}\|^{2}
   -2\eta\langle F(z_{k+\frac12}),z_{k+\frac12}-z^{\star}\rangle
   +\eta^{2}\|F(z_{k+\frac12})\|^{2}\\
&\quad+2\eta^{2}L\|F(z_{k})\|^{2}.
\end{aligned}
\tag{18}
\]

\noindent\textbf{[Step 9]} \emph{Monotonicity of $F$ yields a variational inequality.}  
For any $z\in\mathcal{Z}$,
\[
\langle F(z_{k+\frac12}),z_{k+\frac12}-z\rangle\ge0.
\tag{19}
\]
Choose $z=z^{\star}$ and drop the non‑negative term in \eqref{18}:
\[
\|z_{k+1}-z^{\star}\|^{2}
\le\|z_{k}-z^{\star}\|^{2}
   +\eta^{2}\|F(z_{k+\frac12})\|^{2}
   +2\eta^{2}L\|F(z_{k})\|^{2}.
\tag{20}
\]

\noindent\textbf{[Step 10]} \emph{Bound the gradients by the Lipschitz constant and the diameter.}  
From \eqref{as:smooth} and the boundedness \eqref{as:bounded},
\[
\|F(z)\|\le L\|z-z^{\star}\|\le LD,\qquad\forall z\in\mathcal{Z}.
\tag{21}
\]

\noindent\textbf{[Step 11]} \emph{Insert \eqref{21} into \eqref{20}.}  
\[
\|z_{k+1}-z^{\star}\|^{2}
\le\|z_{k}-z^{\star}\|^{2}
   +\eta^{2}L^{2}D^{2}+2\eta^{2}L^{3}D^{2}
   =\|z_{k}-z^{\star}\|^{2}+3\eta^{2}L^{2}D^{2}.
\tag{22}
\]

\noindent\textbf{[Step 12]} \emph{Telescoping sum.}  
Summing \eqref{22} for $k=0,\dots,T-1$ yields
\[
\|z_{T}-z^{\star}\|^{2}
\le\|z_{0}-z^{\star}\|^{2}+3T\eta^{2}L^{2}D^{2}
\le D^{2}+3T\eta^{2}L^{2}D^{2},
\tag{23}
\]
where we used $\|z_{0}-z^{\star}\|\le D$.

\noindent\textbf{[Step 13]} \emph{Relate the gap to the averaged iterates.}  
For any $z\in\mathcal{Z}$, convexity of $\|\cdot\|^{2}$ gives
\[
\frac{1}{T}\sum_{k=0}^{T-1}\langle F(z_{k+\frac12}),z_{k+\frac12}-z\rangle
\ge\langle F(\bar z_{T}),\bar z_{T}-z\rangle.
\tag{24}
\]
Using the definition of $F$ and the saddle‑point inequality \eqref{3},
\[
\langle F(\bar z_{T}),\bar z_{T}-z\rangle
= L(\bar x_{T},y)-L(\bar x_{T},\bar y_{T})
   +L(\bar x_{T},\bar y_{T})-L(x,\bar y_{T})
   =\gap(\bar z_{T};z).
\tag{25}
\]

\noindent\textbf{[Step 14]} \emph{Upper bound of the left‑hand side of \eqref{24}.}  
From \eqref{8} and \eqref{12} we have for each $k$,
\[
\langle F(z_{k+\frac12}),z_{k+\frac12}-z^{\star}\rangle
\le\frac{1}{2\eta}\bigl(\|z_{k}-z^{\star}\|^{2}
                         -\|z_{k+1}-z^{\star}\|^{2}\bigr).
\tag{26}
\]
Summing over $k$ and dividing by $T$,
\[
\frac{1}{T}\sum_{k=0}^{T-1}\langle F(z_{k+\frac12}),z_{k+\frac12}-z^{\star}\rangle
\le\frac{\|z_{0}-z^{\star}\|^{2}}{2\eta T}
\le\frac{D^{2}}{2\eta T}.
\tag{27}
\]

\noindent\textbf{[Step 15]} \emph{Combine \eqref{24}–\eqref{27}.}  
Using $z=z^{\star}$ in \eqref{24} and the bound \eqref{27},
\[
\gap(\bar z_{T})\le\frac{D^{2}}{2\eta T}.
\tag{28}
\]
Finally, choosing $\eta=1/(2L)$ (the largest stepsize allowed by \eqref{thm:rate}\,(4)) gives
\[
\gap(\bar z_{T})\le\frac{2L D^{2}}{T},
\]
which is exactly the rate claimed in \eqref{5}–\eqref{6}. ∎

\section*{Rate Derivation (With Constants)}
The bound \eqref{28} shows a linear dependence on the diameter $D$ and on the smoothness constant $L$.  
If the feasible set is a Euclidean ball of radius $R$, then $D=2R$ and the rate becomes
\[
\gap(\bar z_{T})\le\frac{8LR^{2}}{T}.
\]

\section*{Special Cases}
\begin{itemize}
\item \textbf{Bilinear games.} $L(x,y)=x^{\top}Ay$ with $A\in\R^{n\times m}$. Here $F(z)=[Ay;\,-A^{\top}x]$ is $L=\|A\|_{2}$‑Lipschitz, and the EG method reduces to the classical “optimistic” gradient method, still enjoying the $O(1/T)$ gap.
\item \textbf{Strongly convex–strongly concave.} If $L$ is $\mu$‑strongly convex in $x$ and $\mu$‑strongly concave in $y$, the same analysis yields a linear convergence rate
\[
\|z_{k}-z^{\star}\|\le\bigl(1-\eta\mu\bigr)^{k}\|z_{0}-z^{\star}\|,
\]
provided $\eta\le 1/L$.
\item \textbf{Stochastic Extra‑Gradient.} When only unbiased stochastic estimates $\tilde F(z)$ are available, the same steps hold in expectation, giving an $O(1/\sqrt{T})$ rate unless variance reduction is employed.
\end{itemize}

\section*{Discussion}
The Extra‑Gradient method achieves the optimal $O(1/T)$ ergodic convergence for smooth convex–concave saddle‑point problems under merely Lipschitz continuity of the gradient and boundedness of the domain. The proof hinges on monotonicity of the associated operator and a clever use of the extrapolation point to cancel the harmful coupling term that plagues plain Gradient Descent–Ascent. The stepsize restriction $\eta\le1/(2L)$ is both sufficient for the analysis and necessary for stability in the worst case. Compared with first‑order baselines, EG enjoys faster decay of the primal–dual gap without any additional computational overhead beyond a second gradient evaluation per iteration.

\end{document}